\appendixitem{STRUCTURE OF CONNECTED COMPACT GROUPS}

Let G be a commutative compact group. Recall (Spectral Theory, Chap. II, §1, no. 9, Prop. 11) that G is then isomorphic to the dual topological group of a discrete commutative group \(\hat{G}\) . The group G is connected if and only if \(\hat{G}\) is torsion-free (Spectral Theory, Chap. II, §2, no. 2, Cor. 1 of Prop. 4).

The following properties are equivalent (Spectral Theory, Chap. II, §2, no. 2, Cor. 2 of Prop. 4 and §1, no. 9, Cor. 2 of Prop. 11):

\begin{enumerate}
\def\labelenumi{(\roman{enumi})}
\item
  G is totally discontinuous;
\item
  \(\hat{\mathrm{G}}\) is a torsion group;
\item
  the topological group G is isomorphic to the limit of a projective system of finite (commutative) groups, each having the discrete topology.
\end{enumerate}

The proposition below generalizes Cor. 1 of Prop. 4 of §1, no. 4.

PROPOSITION 2. Let G be a connected compact group.

\begin{enumerate}
\def\labelenumi{\alph{enumi})}
\item
  \(\mathrm{C}(\mathrm{G})_{0}\) is a commutative connected compact group; \(\mathrm{D}(\mathrm{G})\) is a connected compact group, equal to its derived group.
\item
  The continuous homomorphism \((x,y)\mapsto xy\) from \(\mathrm{C}(\mathrm{G})_{0}\times\mathrm{D}(\mathrm{G})\) to G is surjective and its kernel is a central subgroup of \(\mathrm{C}(\mathrm{G})_{0}\times\mathrm{D}(\mathrm{G})\) that is compact and totally discontinuous.
\item
  There exists a family \((\mathrm{S}_{\lambda})_{\lambda\in\mathrm{L}}\) of almost simple compact Lie groups and a surjective continuous homomorphism \(\prod_{\lambda\in\mathrm{L}}\mathrm{S}_{\lambda}\to\mathrm{D}(\mathrm{G})\) , whose kernel is a totally discontinuous, compact, central subgroup.
\end{enumerate}

Let \((\mathrm{G}_{\alpha}, f_{\alpha\beta})\) be a projective system of compact Lie groups, relative to a filtered set I, such that G is isomorphic to \(\varprojlim G_{\alpha}\) and such that the canonical maps \(f_{\alpha}: G \to G_{\alpha}\) are surjective (Cor. 2 of Prop. 1). For \(\alpha \in I\) , let \(\pi_{\alpha}: \tilde{\mathrm{D}}(\mathrm{G}_{\alpha}) \to \mathrm{D}(\mathrm{G}_{\alpha})\) be a universal covering of the group \(\mathrm{D}(\mathrm{G}_{\alpha})\) . The \(f_{\alpha\beta}\) induce morphisms \(\tilde{f}_{\alpha\beta}: \tilde{\mathrm{D}}(\mathrm{G}_{\beta}) \to \tilde{\mathrm{D}}(\mathrm{G}_{\alpha})\) , and \((\tilde{\mathrm{D}}(\mathrm{G}_{\alpha}), \tilde{f}_{\alpha\beta})\) is a projective system of topological groups satisfying the hypotheses of Lemma 3.

It follows from this lemma that the topological group \(\varprojlim\tilde{\mathrm{D}}(\mathrm{G}_{\alpha})\) is isomorphic to the product of a family \((\mathrm{S}_{\lambda})_{\lambda\in\mathrm{L}}\) of almost simple compact Lie groups. By Lemma 1, the limit of the projective system of homomorphisms \((\pi_{\alpha})\) can be identified with a continuous homomorphism \(\pi:\prod_{\lambda\in\mathrm{L}}\mathrm{S}_{\lambda}\to\overline{\mathrm{D}(\mathrm{G})}\) , which is surjective (General Topology, Chap. I, §9, no. 6, Cor. 2 of Prop. 8).

Now observe that the group
\[
\prod_ {\lambda \in L} S _ {\lambda}
\]

is equal to its derived group: this follows from §4, no. 5, Cor. of Prop. 10. The same is true for \(\overline{\mathrm{D(G)}}\) , since \(\pi\) is surjective. Consequently, \(\mathrm{D(G)} \supset \mathrm{D}(\overline{\mathrm{D(G)}}) = \overline{\mathrm{D(G)}}\) . Thus, the group \(\mathrm{D(G)}\) is compact and equal to its derived group; this proves \(a\) ), since the assertions concerning \(\mathrm{C(G)}_0\) are trivial.

On the other hand, the kernel of \(\pi : \prod_{\lambda \in L} S_{\lambda} \to D(G)\) can be identified with \(\varprojlim \operatorname{Ker}(\pi_{\alpha})\) (Algebra, Chap. II, §6, no. 1, Remark 1), and thus with a compact, totally discontinuous, central subgroup, hence \(c\) ).

We prove \(b\) ). For all \(\alpha\) in I, the morphism \(s_{\alpha}: \mathrm{C}(\mathrm{G}_{\alpha})_0 \times \mathrm{D}(\mathrm{G}_{\alpha}) \to \mathrm{G}_{\alpha}\) such that \(s_{\alpha}(x,y) = xy\) for \(x \in \mathrm{C}(\mathrm{G}_{\alpha})_0, y \in \mathrm{D}(\mathrm{G}_{\alpha})\) , is surjective and its kernel is a finite central subgroup (§1, no. 4, Cor. 1 of Prop. 4). The \(s_{\alpha}\) form a projective system of maps, whose limit can, by the preceding, be identified with the homomorphism \((x,y) \mapsto xy\) from \(\mathrm{C}(\mathrm{G})_0 \times \mathrm{D}(\mathrm{G})\) to \(\mathrm{G}\) . We now see as before that this map is surjective and that its kernel is central and totally discontinuous, hence \(b\) ).

COROLLARY. Every solvable connected compact group is commutative.

Indeed, the derived group is then solvable and equal to its derived group (Prop. 2 a)), hence reduced to the identity element.

